% from ega1-7-fr.tex
\oldpage[I]{162}
\begin{definition}[7.3.2]
\label{I.7.3.2-fr}
On appelle \emph{faisceau des fonctions rationnelles} sur un préschéma $X$
et on désigne par $\mathscr{R}(X)$ la $\mathscr{O}_X$-Algèbre associée au
préfaisceau formé des $R(U)$.
\end{definition}

Pour tout préschéma $X$ et tout ouvert $U\subset X$, il est clair que le
faisceau induit $\mathscr{R}(X)|U$ n'est autre que $\mathscr{R}(U)$.

\begin{proposition}[7.3.3]
\label{I.7.3.3-fr}
Soit $X$ un préschéma tel que la famille $(X_\lambda)$ de ses composantes
irréductibles soit localement finie (ce qui est en particulier le cas lorsque
l'espace sous-jacent à $X$ est localement noethérien). Alors le
$\mathscr{O}_X$-Module $\mathscr{R}(X)$ est quasi-cohérent, et pour tout
ouvert $U$ de $X$ ne rencontrant qu'un nombre fini de composantes $X_\lambda$,
$R(U)$ est égal à $\Gamma(U,\mathscr{R}(X))$ et s'identifie canoniquement au
composé direct des anneaux locaux des points génériques des $X_\lambda$ telles
que $U\cap X_\lambda\neq\varnothing$.
\end{proposition}

On peut évidemment se limiter au cas où $X$ n'a qu'un nombre fini de
composantes irréductibles $X_i$, de points génériques $x_i$
($1\leq i\leq n$). Le fait que $R(U)$ est canoniquement identifié au composé
direct des $\mathscr{O}_{x_i}=R(X_i)$ tels que
$U\cap X_i\neq\varnothing$ résulte alors de
(\hyperref[I.7.1.7-fr]{7.1.7}). Montrons en outre que le préfaisceau
$U\to R(U)$ vérifie les axiomes des faisceaux, ce qui prouvera que
$R(U)=\Gamma(U,\mathscr{R}(X))$. En effet, il vérifie (F 1) d'après ce qui
précède. Pour voir qu'il satisfait à (F 2), considérons un recouvrement d'un
ouvert $U$ de $X$ par des ouverts $V_\alpha\subset U$~; si les
$s_\alpha\in R(V_\alpha)$ sont telles que les restrictions de $s_\alpha$ et
$s_\beta$ à $V_\alpha\cap V_\beta$ coïncident pour tout couple d'indices, on
en conclut que pour tout indice $i$ tel que $U\cap X_i\neq\varnothing$, les
composantes dans $R(X_i)$ de toutes les $s_\alpha$ telles que
$V_\alpha\cap X_i\neq\varnothing$ sont les mêmes~; désignant par $t_i$ cette
composante, il est clair que l'élément de $R(U)$ ayant les $t_i$ pour
composantes a pour restriction $s_\alpha$ à chaque $V_\alpha$. Enfin, pour
voir que $\mathscr{R}(X)$ est quasi-cohérent, on peut se limiter au cas où
$X=\operatorname{Spec}(A)$ est affine~; en prenant pour $U$ les ouverts
affines de la forme $D(f)$, où $f\in A$, il résulte de ce qui précède et de la
définition (\hyperref[I.1.3.4-fr]{1.3.4}) que l'on a
$\mathscr{R}(X)=\widetilde M$, où $M$ est somme directe des $A$-modules
$A_{x_i}$.

\begin{corollary}[7.3.4]
\label{I.7.3.4-fr}
Soit $X$ un préschéma réduit n'ayant qu'un nombre fini de composantes
irréductibles, et soient $X_i$ ($1\leq i\leq n$) les sous-préschémas fermés
réduits de $X$ ayant pour espaces sous-jacents les composantes irréductibles de
$X$ (\hyperref[I.5.2.1-fr]{5.2.1}). Si $h_i$ est l'injection canonique
$X_i\to X$, $\mathscr{R}(X)$ est alors composée directe des
$\mathscr{O}_X$-Algèbres $(h_i)_*(\mathscr{R}(X_i))$.
\end{corollary}

\begin{corollary}[7.3.5]
\label{I.7.3.5-fr}
Si $X$ est irréductible, tout $\mathscr{R}(X)$-Module quasi-cohérent
$\mathscr{F}$ est un faisceau simple.
\end{corollary}

Il suffit de montrer que tout $x\in X$ admet un voisinage $U$ tel que
$\mathscr{F}|U$ soit un faisceau simple
(\hyperref[0.3.6.2-fr]{0, 3.6.2}), autrement dit on est ramené au cas où $X$
est affine~; on peut en outre supposer que $\mathscr{F}$ est le conoyau d'un
homomorphisme
$(\mathscr{R}(X))^{(I)}\to(\mathscr{R}(X))^{(J)}$
(\hyperref[0.5.1.3-fr]{0, 5.1.3}), et tout revient à voir que
$\mathscr{R}(X)$ est un faisceau simple~; mais cela est évident puisque
$\Gamma(U,\mathscr{R}(X))=R(X)$ pour tout ouvert non vide $U$, $U$ contenant
le point générique de $X$.

\begin{corollary}[7.3.6]
\label{I.7.3.6-fr}
Si $X$ est irréductible, pour tout $\mathscr{O}_X$-Module quasi-cohérent
$\mathscr{F}$, $\mathscr{F}\otimes_{\mathscr{O}_X}\mathscr{R}(X)$ est un
faisceau simple~; si en outre $X$ est réduit (donc intègre),
$\mathscr{F}\otimes_{\mathscr{O}_X}\mathscr{R}(X)$ est isomorphe à un
faisceau de la forme $(\mathscr{R}(X))^{(I)}$.
\end{corollary}

La seconde assertion résulte de ce que $R(X)$ est alors un corps.

\begin{proposition}[7.3.7]
\label{I.7.3.7-fr}
Supposons que le préschéma $X$ soit localement intègre ou localement
